\section{Seals}\label{sec:seals}

Back at the committee, Sita's proposal has its endorsements and the vote passes. The committee's
rule says a decision is final only when no objection remains open. The chair looks around the room
and asks, ``Any objections?'' There is a pause. ``Hearing none, the motion carries.'' The ritual is
older than any computer, and it does what the CALM theorem says must be done. The chair cannot know
that no member, anywhere, has an objection in their pocket. What the chair can do is declare that
the objections raised by now are the ones that count. After that, a member who objects next week
starts a new proposal, and this one stays closed.

\citet{smithsearle2003} pressed Searle with cases of contested property, where ownership rests on
what Smith called a negative basis, the absence of any competing claim, and Searle answered that
problems of that kind are addressed by legal or political means. The sort says why both are right. A
court closes the record: it declares that the claims before it are all the claims that count.

We call such a declaration a \emph{seal}. A land registry seals the titles to a plot, and an editor
closes an issue. Each performs a status function declaration in Searle's sense, making it the case
that a set of records counts as complete for a question. Once the seal exists, the absence fact can
be grounded in the seal itself, which is a record like any other. The decision of the meeting holds
as long as the seal exists, and adding records can never remove it, so the fact has become monotone
again. Coordination was spent once, at the moment of sealing, and after that the fact travels freely.
Seen this way, an authority may know no more than anyone else. In a world of records kept in many
hands nobody can know that nothing is missing, and what an authority has is the standing to make
what is missing stop counting.

\subsection{Sealing together}\label{sec:joint}

A seal can come from somewhere other than a chair, a court or a registry. In Gilbert's best-known
example, two friends agree to walk to the market together, and halfway there one of them speeds up
and leaves the other behind without a word \citep{gilbert1990}. Something has gone wrong, and both
know it. The one left behind has the standing to object, and neither of them can call the walk off
alone: it took both to begin it, and it takes both to end it. Gilbert calls what holds them a joint
commitment, through which two or more people form a plural subject with commitments of its own that
no single member can undo \citep{gilbert2013}.

A seal can be made this way. Every party the question names signs a record naming the same set of
records, and the seal exists when all the signatures do. Each signature is an open and checkable
expression of readiness, every signer can see who else signed, and none of them can unseal it alone:
reopening takes another record signed by the same parties, so even a change of mind becomes one more
record in the history. At the clinic this is a handover. Before the day doctor leaves, she and the
night nurse co-sign the set of prescriptions the current dose is computed from, and the patient can
sign it too. The nurse no longer has to guess whether a newer prescription is hiding somewhere,
since she works from a set that the people responsible for it sealed together. The authority to seal
remains, shared among the people the question concerns.

\subsection{What a signature can vouch for}\label{sec:vouch}

Cameras have begun to sign their photographs, attaching a credential that says which device made a
picture and that nothing has changed since. In September 2025 a photographer, Adam Horshack, showed
that a camera with this feature could be made to sign a picture it never took, by overlaying an
image made elsewhere in its multiple-exposure mode, and the result passed the public verification
tool \citep{petapixel2025a}. The manufacturer suspended the service and later announced that it would
revoke every certificate it had issued \citep{petapixel2025b}.

The case shows what a signature can make true. A content credential
checks that a record is well formed, that it is untouched, and that a signer on a trusted list signed
it \citep{c2pa}. Those checks establish the kind of fact a signature can ground: this signer vouched
for this file. Whether the scene in the file took place is a different kind of fact, and no signature
settles it. As generative models write more records of events that never happened, the question to
ask of a record moves from what it shows to who stands behind it.

It also shows where absence enters. A credential stays valid only while its certificate has not been
revoked. Statuses of this kind are defeasible in the lawyer's sense \citep{hart1949}: they hold
unless something defeats them, such as a revocation, an appeal, or an exception that comes to light
later. Every defeasible status, from a valid license to a current dose, hides an absence. A verifier
that never received the revocation goes on calling the credential valid, with nobody having done
anything wrong, because its grounding condition asked for an absence and nobody had sealed the list
of revocations it consulted. A revocation list that is published, signed and dated is such a seal,
and a verifier that checks one can say that a credential was valid as of that list.

\subsection{Where authority goes}\label{sec:where}

The systems that now hold our records each settle facts in their own way. A platform that holds
everyone's records settles all of them, monotone or not, and becomes the authority over facts that
never needed one. A global ledger makes every participant agree on every fact, including the many
that only ask for a record to exist. The sort points to a third design. Records stay with the people
who signed them, every monotone fact is settled by whoever holds the records, and coordination is
spent only where a fact asks for an absence, among the people that fact concerns, who seal it
together. Authority is needed where absence matters, and where it is needed, it can be shared. That
leaves the question of who may hold it.
